% GENERATED FILE. Edit canonical lessons, then run npm run generate:books.
% canonical-source-hash: fnv1a64:98ead500ea672a43
% canonical-lessons: MW-C98-jimano, MW-C98-chauk, MW-C98-chauraho, MW-C98-mor, MW-C98-gelo

\chapter{Right, Town square, Crossroads, Turn, Path}
\label{ch:mw-right-town-square-crossroads-turn-path}
\hlchaptermodality{98}

\begin{chapteropening}
\textbf{By the end of this chapter:} \emph{I can say right, a town square, a crossroads, a turn and a path.}

Complete the last lesson of chapter 98: I can say right, a town square, a crossroads, a turn and a path.
\end{chapteropening}

\section[\texorpdfstring{\mw{जीमणो} (jīmaṇo)}{जीमणो (jīmaṇo)}]{\mw{जीमणो} (jīmaṇo) — right (the side)}

\label{lesson:MW-C98-jimano}



\begin{quote}
Before the new one: say the Marwadi for opposite, then the Marwadi for left.
\end{quote}

\subsection*{You'll want to know: \mw{जीमणो}}
\textbf{\mw{जीमणो}} — \emph{jīmaṇo} — \textquotedblleft{}right (the side)\textquotedblright{}.

\begin{cousinweb}
\textbf{\mw{जीमणो}} (\emph{jīmaṇo}) is a tadbhava, a word inherited from Sanskrit and reshaped by sound change. It comes from \emph{jemana}, \textquotedblleft{}eating\textquotedblright{}. CDIAL lists Old Marwari \emph{jīmaṇa} under it (CDIAL 5268). The Old Marwari row means \textquotedblleft{}food, meal\textquotedblright{}, and Old Gujarati \emph{jimaṇaüṃ} means \textquotedblleft{}to the righthand side\textquotedblright{}: the right is the eating hand.
\end{cousinweb}

\begin{grammarlens}[title={What you've built}]
One more word for this chapter.
\end{grammarlens}

\subsection*{Your turn}
\begin{itemize}
\raggedright
  \item \emph{Say it:} \emph{jīmaṇo}
  \item \emph{Say it:} \emph{jīmaṇo}, once more
  \item \emph{Say it:} say \emph{jīmaṇo}
  \item \emph{Recall:} say the Marwadi for opposite, then the Marwadi for left, then say \emph{jīmaṇo} again
\end{itemize}

\begin{tcolorbox}[breakable,colback=teal!4,colframe=teal!35!black,title={Before you move on}]
What does \mw{जीमणो} mean? (\textquotedblleft{}Right (the side)\textquotedblright{}.) Say it once more.
\end{tcolorbox}
\section[\texorpdfstring{\mw{चौक} (chauk)}{चौक (chauk)}]{\mw{चौक} (chauk) — a town square}

\label{lesson:MW-C98-chauk}



\begin{quote}
Before the new one: say the Marwadi for left, then the Marwadi for right.
\end{quote}

\subsection*{You'll want to know: \mw{चौक}}
\textbf{\mw{चौक}} — \emph{chauk} — \textquotedblleft{}a town square\textquotedblright{}.

\begin{cousinweb}
\textbf{\mw{चौक}} (\emph{chauk}) is a tadbhava, a word inherited from Sanskrit and reshaped by sound change. It comes from \emph{catuṣka}, \textquotedblleft{}set of four; square courtyard\textquotedblright{}. CDIAL gives no Marwari form in this entry (CDIAL 4629); the match is through Hindi \emph{cauk}. It is built on \emph{catur}, \textquotedblleft{}four\textquotedblright{}.
\end{cousinweb}

\begin{grammarlens}[title={What you've built}]
One more word for this chapter.
\end{grammarlens}

\subsection*{Your turn}
\begin{itemize}
\raggedright
  \item \emph{Say it:} \emph{chauk}
  \item \emph{Say it:} \emph{chauk}, once more
  \item \emph{Say it:} say \emph{chauk}
  \item \emph{Recall:} say the Marwadi for left, then the Marwadi for right, then say \emph{chauk} again
\end{itemize}

\begin{tcolorbox}[breakable,colback=teal!4,colframe=teal!35!black,title={Before you move on}]
What does \mw{चौक} mean? (\textquotedblleft{}A town square\textquotedblright{}.) Say it once more.
\end{tcolorbox}
\section[\texorpdfstring{\mw{चौराहो} (chaurāho)}{चौराहो (chaurāho)}]{\mw{चौराहो} (chaurāho) — a crossroads}

\label{lesson:MW-C98-chauraho}



\begin{quote}
Before the new one: say the Marwadi for right, then the Marwadi for a town square.
\end{quote}

\subsection*{You'll want to know: \mw{चौराहो}}
\textbf{\mw{चौराहो}} — \emph{chaurāho} — \textquotedblleft{}a crossroads\textquotedblright{}.

\begin{cousinweb}
\textbf{\mw{चौराहो}} (\emph{chaurāho}) is a meeting of four roads: \emph{chau-}, \textquotedblleft{}four\textquotedblright{}, as in \emph{chauk} (CDIAL 4629), with \emph{rāh}, \textquotedblleft{}road\textquotedblright{}, usually traced to Persian \emph{rāh}, \textquotedblleft{}road, way\textquotedblright{}. CDIAL has no entry for the compound.
\end{cousinweb}

\begin{grammarlens}[title={What you've built}]
One more word for this chapter.
\end{grammarlens}

\subsection*{Your turn}
\begin{itemize}
\raggedright
  \item \emph{Say it:} \emph{chaurāho}
  \item \emph{Say it:} \emph{chaurāho}, once more
  \item \emph{Say it:} say \emph{chaurāho}
  \item \emph{Recall:} say the Marwadi for right, then the Marwadi for a town square, then say \emph{chaurāho} again
\end{itemize}

\begin{tcolorbox}[breakable,colback=teal!4,colframe=teal!35!black,title={Before you move on}]
What does \mw{चौराहो} mean? (\textquotedblleft{}A crossroads\textquotedblright{}.) Say it once more.
\end{tcolorbox}
\section[\texorpdfstring{\mw{मोड़} (moṛ)}{मोड़ (moṛ)}]{\mw{मोड़} (moṛ) — a turn, a bend}

\label{lesson:MW-C98-mor}



\begin{quote}
Before the new one: say the Marwadi for a town square, then the Marwadi for a crossroads.
\end{quote}

\subsection*{You'll want to know: \mw{मोड़}}
\textbf{\mw{मोड़}} — \emph{moṛ} — \textquotedblleft{}a turn, a bend\textquotedblright{}.

\begin{cousinweb}
\textbf{\mw{मोड़}} (\emph{moṛ}) is a tadbhava, a word inherited from Sanskrit and reshaped by sound change. It comes from \emph{muṭati}, \textquotedblleft{}twists\textquotedblright{}. CDIAL gives no Marwari form in this entry (CDIAL 10186); the match is through Hindi \emph{moṛnā} and Gujarati \emph{moṛvũ}. It is the noun of the verb \textquotedblleft{}to turn, to bend\textquotedblright{}. CDIAL notes that the Sanskrit verb partly ran together with a look-alike it takes from Dravidian.
\end{cousinweb}

\begin{grammarlens}[title={What you've built}]
One more word for this chapter.
\end{grammarlens}

\subsection*{Your turn}
\begin{itemize}
\raggedright
  \item \emph{Say it:} \emph{moṛ}
  \item \emph{Say it:} \emph{moṛ}, once more
  \item \emph{Say it:} say \emph{moṛ}
  \item \emph{Recall:} say the Marwadi for a town square, then the Marwadi for a crossroads, then say \emph{moṛ} again
\end{itemize}

\begin{tcolorbox}[breakable,colback=teal!4,colframe=teal!35!black,title={Before you move on}]
What does \mw{मोड़} mean? (\textquotedblleft{}A turn, a bend\textquotedblright{}.) Say it once more.
\end{tcolorbox}
\section[\texorpdfstring{\mw{गेलो} (gelo)}{गेलो (gelo)}]{\mw{गेलो} (gelo) — a path, the way}

\label{lesson:MW-C98-gelo}



\begin{quote}
Before the new one: say the Marwadi for a crossroads, then the Marwadi for a turn.
\end{quote}

\subsection*{You'll want to know: \mw{गेलो}}
\textbf{\mw{गेलो}} — \emph{gelo} — \textquotedblleft{}a path, the way\textquotedblright{}.

\begin{cousinweb}
\textbf{\mw{गेलो}} (\emph{gelo}) is a tadbhava, a word inherited from Sanskrit and reshaped by sound change. It comes from \emph{gati}, \textquotedblleft{}going, gait\textquotedblright{}. CDIAL lists Old Marwari \emph{gailo} under it (CDIAL 4009). The Old Marwari row means \textquotedblleft{}road\textquotedblright{}; Hindi has \emph{gail}, \textquotedblleft{}path\textquotedblright{}.
\end{cousinweb}

\begin{grammarlens}[title={What you've built}]
One more word for this chapter.
\end{grammarlens}

\subsection*{Your turn}
\begin{itemize}
\raggedright
  \item \emph{Say it:} \emph{gelo}
  \item \emph{Say it:} \emph{gelo}, once more
  \item \emph{Say it:} say \emph{gelo}
  \item \emph{Recall:} say the Marwadi for a crossroads, then the Marwadi for a turn, then say \emph{gelo} again
\end{itemize}

\begin{tcolorbox}[breakable,colback=teal!4,colframe=teal!35!black,title={Before you move on}]
What does \mw{गेलो} mean? (\textquotedblleft{}A path, the way\textquotedblright{}.) Say it once more.
\end{tcolorbox}
